% Generated by publication/build_sections.py; do not edit by hand.
\paragraph{Informal verdict.}
\textbf{A --- Complete and apparently correct}, with \textbf{high confidence}.

\paragraph{Lean coverage.}
Full canonical Q865 coverage. The preserved 1,653-line Q865.lean gives the global product definition and natural-evaluation formula for every A\_k, summability and the exact positive-root series for every n >= 2, and the odd-index limit for rmin. Q865Completion.lean proves polynomial nonzeroness and root existence for every m >= 2, finiteness/nonemptiness of the root-modulus set, attainment of its infimum, universal minimality, an IsLeast form, and the unconditional Q865\_answer synthesis theorem. The two modules are not yet imported into lean/RMS.

\subsection{Audited informal answer}
Q865 --- Complete solution
\begin{enumerate}
\item Verified statement
\end{enumerate}
For \(n\ge 2\), define

\[
P_n(X)=X^n-\sum_{k=0}^{n-1}X^k.
\]

Let \(\rho_n\) be its unique positive real root, let \(r_n\) be the smallest modulus of its complex roots, and put

\[
u_n=2^{-n-1}.
\]

The earlier response R469 established an asymptotic expansion

\[
\rho_n
=
2\left(
1-\sum_{k=0}^{N}A_k(n)u_n^{k+1}
+o(u_n^{N+1})
\right),
\]

where the \(A_k\) are polynomials, together with \(r_n\to1\) and

\[
2n(1-r_{2n})\longrightarrow \log 3.
\]

Q865 asks for a closed formula for \(A_k\), whether the series for \(\rho_n\) converges exactly, and whether the analogous odd-index limit equals \(\log 3\). There is no apparent ambiguity or typo in the displayed statement. Luc Pommeret
\begin{enumerate}
\item Final theorem
\end{enumerate}
For every \(k\ge0\),

\[
\boxed{
A_k(X)
=
\frac1{k+1}
\binom{(k+1)(X+1)-2}{k}
=
\frac1{(k+1)!}
\prod_{j=0}^{k-1}\bigl((k+1)X+j\bigr)
}
\]

with the empty product interpreted as \(1\). Thus, for example,

\[
A_0(X)=1,\qquad
A_1(X)=X,\qquad
A_2(X)=\frac{X(3X+1)}2.
\]

For every integer \(n\ge2\), the series in part b converges absolutely and

\[
\boxed{
\rho_n
=
2\left(
1-\sum_{k=0}^{\infty}A_k(n)u_n^{k+1}
\right).
}
\]

Equivalently,

\[
\boxed{
\rho_n
=
2\left(
1-
\sum_{m=1}^{\infty}
\frac1m
\binom{(n+1)m-2}{m-1}
\,2^{-m(n+1)}
\right).
}
\]

Finally, a stronger statement than part c holds:

\[
\boxed{
r_m
=
1-\frac{\log 3}{m}
+
\frac{\frac12(\log3)^2+\frac13\log3}{m^2}
+
O(m^{-3})
\qquad(m\to\infty).
}
\]

In particular,

\[
\boxed{
\lim_{n\to\infty}(2n+1)(1-r_{2n+1})=\log3.
}
\]


\begin{enumerate}
\item Algebraic reduction
\end{enumerate}
A direct calculation gives

\[
(z-1)P_n(z)
=
z^{n+1}-2z^n+1
=
1-z^n(2-z).
\]

Consequently, for \(z\ne1\),

\[
P_n(z)=0
\quad\Longleftrightarrow\quad
z^n(2-z)=1.
\tag{3.1}
\]

For the positive root, write

\[
\rho_n=2(1-y_n).
\]

Then (3.1) becomes

\[
2^n(1-y_n)^n\,2y_n=1,
\]

or

\[
y_n(1-y_n)^n=2^{-n-1}=u_n.
\tag{3.2}
\]

The relevant solution satisfies

\[
0<y_n<\frac1{n+1}.
\tag{3.3}
\]

Indeed, the function

\[
F_n(y)=y(1-y)^n
\]

is strictly increasing on \([0,1/(n+1)]\), with maximum there equal to

\[
\mathcal R_n
=
F_n\left(\frac1{n+1}\right)
=
\frac{n^n}{(n+1)^{n+1}}.
\tag{3.4}
\]


\begin{enumerate}
\item Inverting \(y(1-y)^n=u\)
\end{enumerate}
We first establish the coefficient formula directly.
Lemma
Let \(n\ge1\) be fixed. The analytic solution near \(u=0\) of

\[
Y=u(1-Y)^{-n},
\qquad Y(0)=0,
\tag{4.1}
\]

has the expansion

\[
Y(u)
=
\sum_{m=1}^{\infty}c_m(n)u^m,
\]

where

\[
c_m(n)
=
\frac1m
\binom{(n+1)m-2}{m-1}.
\tag{4.2}
\]

Proof
Put

\[
\Phi(t)=(1-t)^{-n}.
\]

Equation (4.1) is \(Y=u\Phi(Y)\). For \(m\ge1\), Cauchy's coefficient formula gives

\[
[u^m]Y(u)
=
\operatorname*{Res}_{u=0}
\frac{Y(u)}{u^{m+1}}\,du.
\]

Using the local substitution

\[
u=\frac{t}{\Phi(t)},
\]

we obtain

\[
\begin{aligned}
[u^m]Y(u)
&=
\operatorname*{Res}_{t=0}
\left(
t^{-m}\Phi(t)^m
-
t^{-m+1}\Phi(t)^{m-1}\Phi'(t)
\right)dt.
\end{aligned}
\]

Now

\[
\Phi^{m-1}\Phi'
=
\frac1m(\Phi^m)'.
\]

Moreover, the residue of a derivative vanishes, so

\[
0
=
\operatorname*{Res}_{t=0}
\left(t^{-m+1}\Phi(t)^m\right)'
=
-(m-1)\operatorname*{Res}_{t=0}t^{-m}\Phi(t)^m
+
\operatorname*{Res}_{t=0}t^{-m+1}(\Phi(t)^m)'.
\]

It follows that

\[
[u^m]Y(u)
=
\frac1m
\operatorname*{Res}_{t=0}t^{-m}\Phi(t)^m
=
\frac1m[t^{m-1}]\Phi(t)^m.
\]

Since

\[
\Phi(t)^m=(1-t)^{-nm},
\]

we find

\[
[t^{m-1}](1-t)^{-nm}
=
\binom{nm+m-2}{m-1},
\]

which proves (4.2). \(\square\)

\begin{enumerate}
\item Answer to part a
\end{enumerate}
Applying the lemma to (3.2), we have formally

\[
y_n
=
\sum_{m\ge1}
\frac1m
\binom{(n+1)m-2}{m-1}u_n^m.
\]

Writing \(m=k+1\), the coefficient of \(u_n^{k+1}\) is therefore

\[
A_k(n)
=
\frac1{k+1}
\binom{(n+1)(k+1)-2}{k}.
\]

Since

\[
(n+1)(k+1)-2=(k+1)n+k-1,
\]

this becomes

\[
A_k(n)
=
\frac1{k+1}
\binom{(k+1)n+k-1}{k}.
\]

As a polynomial in an arbitrary variable \(X\),

\[
\boxed{
A_k(X)
=
\frac1{k+1}\binom{(k+1)X+k-1}{k}.
}
\]

Expanding the binomial polynomial,

\[
\binom{(k+1)X+k-1}{k}
=
\frac1{k!}
\prod_{j=0}^{k-1}\bigl((k+1)X+j\bigr),
\]

and hence

\[
\boxed{
A_k(X)
=
\frac1{(k+1)!}
\prod_{j=0}^{k-1}\bigl((k+1)X+j\bigr).
}
\]

This is indeed a polynomial of degree \(k\).

\begin{enumerate}
\item Answer to part b: convergence and exact equality
\end{enumerate}
6.1 Radius of convergence
For fixed \(n\),

\[
c_m(n)
=
\frac{((n+1)m-2)!}{m!(nm-1)!}.
\]

Using

\[
\log(M!)=M\log M-M+O(\log M),
\]

we obtain

\[
\lim_{m\to\infty}c_m(n)^{1/m}
=
\frac{(n+1)^{n+1}}{n^n}.
\]

Thus the radius of convergence of the inverse series is exactly

\[
\boxed{
\mathcal R_n=\frac{n^n}{(n+1)^{n+1}}.
}
\tag{6.1}
\]

This agrees with the critical value (3.4).
We now check that \(u_n\) lies strictly inside that disk. Since \(n/(n+1)\ge2/3\),

\[
\mathcal R_n
=
\frac1{n+1}\left(\frac{n}{n+1}\right)^n
\ge
\frac1{n+1}\left(\frac23\right)^n.
\]

For \(n\ge2\),

\[
\frac{(2/3)^n/(n+1)}{2^{-n-1}}
=
\frac{2(4/3)^n}{n+1}>1.
\]

Indeed, the latter expression equals \(32/27\) at \(n=2\), and its ratio at consecutive indices is

\[
\frac{4(n+1)}{3(n+2)}\ge1
\qquad(n\ge2).
\]

Therefore

\[
u_n<\mathcal R_n
\qquad(n\ge2).
\tag{6.2}
\]

6.2 Identification of the sum
Let \(Y_n(u)\) denote the sum of the power series. The polynomial identity

\[
Y_n(u)(1-Y_n(u))^n=u
\tag{6.3}
\]

holds near \(0\), hence throughout the disk \(|u|<\mathcal R_n\) by the identity theorem.
For \(0<u<\mathcal R_n\), all coefficients of \(Y_n\) are positive. Moreover, \(Y_n(u)\) cannot reach \(1/(n+1)\), since that would give, by (6.3),

\[
u=\frac{n^n}{(n+1)^{n+1}}=\mathcal R_n.
\]

Thus

\[
0<Y_n(u)<\frac1{n+1}
\qquad(0<u<\mathcal R_n).
\]

At \(u=u_n\), set

\[
\widetilde\rho_n=2(1-Y_n(u_n)).
\]

Then \(\widetilde\rho_n>2n/(n+1)>1\), and (6.3) gives

\[
\widetilde\rho_n^n(2-\widetilde\rho_n)=1.
\]

Therefore \(P_n(\widetilde\rho_n)=0\). It is the unique positive root: on \((1,2)\), the function

\[
x\longmapsto x^n(2-x)
\]

increases up to \(2n/(n+1)\), then decreases strictly to \(0\); besides the extraneous solution \(x=1\), the equation \(x^n(2-x)=1\) consequently has exactly one solution.
Hence \(\widetilde\rho_n=\rho_n\), proving

\[
\boxed{
\rho_n
=
2\left(
1-\sum_{k=0}^{\infty}A_k(n)u_n^{k+1}
\right).
}
\]

6.3 Recovery of the stated asymptotic expansion
For completeness, the exact series implies the required remainder estimate uniformly for each fixed truncation order.
For \(m\ge1\),

\[
c_m(n)
\le
\binom{(n+1)m}{m}
\le
\bigl(e(n+1)\bigr)^m.
\]

Let

\[
q_n=e(n+1)u_n.
\]

Then \(q_n\to0\). For every fixed \(N\), and all sufficiently large \(n\),

\[
\begin{aligned}
0
&\le
Y_n(u_n)-\sum_{m=1}^{N+1}c_m(n)u_n^m\\
&\le
\sum_{m=N+2}^{\infty}q_n^m
=
\frac{q_n^{N+2}}{1-q_n}.
\end{aligned}
\]

Consequently,

\[
Y_n(u_n)-\sum_{m=1}^{N+1}c_m(n)u_n^m
=
O_N\!\left((n+1)^{N+2}u_n^{N+2}\right)
=
o(u_n^{N+1}),
\]

because \((n+1)^{N+2}u_n\to0\).

\begin{enumerate}
\item Answer to part c
\end{enumerate}
To avoid conflicting indices, write \(m\) for the degree and retain \(r_m\) for the smallest root modulus.
7.1 A universal lower bound
Let \(s_m\in(0,1)\) be the unique solution of

\[
s_m^m(2+s_m)=1.
\tag{7.1}
\]

Uniqueness follows because \(s\mapsto s^m(2+s)\) is strictly increasing on \((0,\infty)\).
If \(z\) is any root of \(P_m\), then by (3.1),

\[
1=|z|^m|2-z|
\le
|z|^m(2+|z|).
\]

By the definition of \(s_m\),

\[
|z|\ge s_m.
\]

Therefore

\[
r_m\ge s_m.
\tag{7.2}
\]

For even \(m\), equality holds: the number \(-s_m\) satisfies

\[
(-s_m)^m(2+s_m)=1,
\]

and hence is a root of \(P_m\). Thus

\[
r_m=s_m
\qquad(m\ \text{even}).
\tag{7.3}
\]

7.2 A near-negative root for odd \(m\)
Suppose now that

\[
m=2q+1\ge3.
\]

Put

\[
\theta_0=\pi-\frac{\pi}{m}.
\]

For \(\theta\in[\theta_0,\pi]\), consider

\[
H_\theta(t)=t^m|2-te^{i\theta}|,
\qquad 0\le t\le1.
\]

On this interval, \(\cos\theta<0\), and

\[
\frac{\partial}{\partial t}|2-te^{i\theta}|^2
=
2t-4\cos\theta>0.
\]

Thus \(H_\theta\) is strictly increasing. Since

\[
H_\theta(0)=0,
\qquad
H_\theta(1)=|2-e^{i\theta}|>1,
\]

there is a unique \(t(\theta)\in(0,1)\) such that

\[
t(\theta)^m|2-t(\theta)e^{i\theta}|=1.
\tag{7.4}
\]

The function \(t(\theta)\) is continuous.
Let

\[
\alpha(\theta)
=
\operatorname{Arg}\bigl(2-t(\theta)e^{i\theta}\bigr),
\]

where the argument is chosen in \((-\pi/2,0]\). It is continuous, negative for \(\theta<\pi\), and \(\alpha(\pi)=0\). Define

\[
\Phi(\theta)=m\theta+\alpha(\theta).
\]

At the left endpoint,

\[
m\theta_0=(m-1)\pi=2q\pi,
\]

so

\[
\Phi(\theta_0)<2q\pi.
\]

At the right endpoint,

\[
\Phi(\pi)=m\pi=(2q+1)\pi>2q\pi.
\]

The intermediate value theorem therefore gives some

\[
\theta_m\in\left(\pi-\frac{\pi}{m},\pi\right)
\]

such that

\[
m\theta_m+\alpha(\theta_m)=2q\pi.
\]

Let

\[
t_m=t(\theta_m),
\qquad
z_m=t_me^{i\theta_m}.
\]

Equation (7.4) gives the modulus of \(z_m^m(2-z_m)\), while the last phase equation shows that its argument is \(2q\pi\). Hence

\[
z_m^m(2-z_m)=1.
\]

Since \(z_m\ne1\), this gives \(P_m(z_m)=0\). Therefore

\[
s_m\le r_m\le t_m<1.
\tag{7.5}
\]

7.3 Asymptotics of \(s_m\)
Let

\[
L=\log3,
\qquad
x_m=1-s_m.
\]

From (7.1),

\[
s_m^m=\frac1{2+s_m}>\frac13,
\]

so

\[
s_m>3^{-1/m},
\qquad
x_m=O(m^{-1}).
\]

Equation (7.1) becomes

\[
-m\log(1-x_m)=\log(3-x_m).
\]

Using

\[
-\log(1-x)=x+\frac{x^2}{2}+O(x^3),
\]

and

\[
\log(3-x)=L-\frac{x}{3}+O(x^2),
\]

we obtain

\[
m x_m+\frac m2x_m^2+O(mx_m^3)
=
L-\frac{x_m}{3}+O(x_m^2).
\]

First this gives \(x_m=L/m+O(m^{-2})\). Substituting this estimate back yields

\[
x_m
=
\frac{L}{m}
-
\frac{\frac12L^2+\frac13L}{m^2}
+
O(m^{-3}).
\]

Thus

\[
s_m
=
1-\frac{L}{m}
+
\frac{\frac12L^2+\frac13L}{m^2}
+
O(m^{-3}).
\tag{7.6}
\]

7.4 Comparison of \(t_m\) and \(s_m\)
For odd \(m\), write

\[
\theta_m=\pi-\delta_m,
\qquad
0<\delta_m<\frac{\pi}{m}.
\]

Then

\[
|2-t_me^{i\theta_m}|
=
|2+t_me^{-i\delta_m}|,
\]

and

\[
|2+t_me^{-i\delta_m}|^2
=
(2+t_m)^2-4t_m(1-\cos\delta_m).
\]

Therefore

\[
\log(2+t_m)
-
\log|2-t_me^{i\theta_m}|
=
O(\delta_m^2)
=
O(m^{-2}).
\tag{7.7}
\]

Define

\[
G_m(x)=-m\log x-\log(2+x).
\]

We have \(G_m(s_m)=0\). Since \(t_m\) satisfies

\[
-m\log t_m
=
\log|2-t_me^{i\theta_m}|,
\]

equation (7.7) gives

\[
G_m(t_m)=O(m^{-2}).
\]

Moreover,

\[
|G_m'(x)|
=
\frac{m}{x}+\frac1{2+x}
\ge m
\qquad(0<x\le1).
\]

The mean value theorem now gives

\[
t_m-s_m=O(m^{-3}).
\tag{7.8}
\]

Combining (7.5), (7.6), and (7.8), we conclude that for odd \(m\),

\[
r_m
=
1-\frac{L}{m}
+
\frac{\frac12L^2+\frac13L}{m^2}
+
O(m^{-3}).
\]

For even \(m\), the same formula follows from \(r_m=s_m\). Hence it holds for all \(m\):

\[
\boxed{
r_m
=
1-\frac{\log3}{m}
+
\frac{\frac12(\log3)^2+\frac13\log3}{m^2}
+
O(m^{-3}).
}
\]

Taking \(m=2n+1\) gives

\[
(2n+1)(1-r_{2n+1})
=
\log3
-
\frac{\frac12(\log3)^2+\frac13\log3}{2n+1}
+
O(n^{-2}),
\]

and therefore

\[
\boxed{
\lim_{n\to\infty}(2n+1)(1-r_{2n+1})=\log3.
}
\]

\begin{enumerate}
\item Boundary checks
\end{enumerate}
For \(n=2\),

\[
P_2(X)=X^2-X-1,
\qquad
\rho_2=\frac{1+\sqrt5}{2},
\qquad
r_2=\frac{\sqrt5-1}{2}.
\]

The latter is exactly the solution of

\[
s^2(2+s)=1,
\]

as predicted by the even-degree argument.
The restriction \(n\ge2\) also avoids the degenerate case \(n=1\): then \(P_1(X)=X-1\), while

\[
u_1=\frac14=\mathcal R_1
\]

lies exactly at the branch point of the inverse \(y(1-y)=u\), rather than strictly inside its disk of convergence.
